%% kapa: exam=2 round=35 year=2024 period=3 subject=compensation_law subject_ko=감정평가 및 보상법규 seq=1 number=1 points=40
%% source: Q-Net 감정평가사 기출문제 (cst00309) artlSeq=5213532 file=raw/1008/5213532/2024년 제35회 감정평가사 2차시험 문제지.pdf
%% conversion: kordoc 4.24.51 -> tools/gen2.py (auto); review=proofread-sonnet; flags=none
\begin{longquestion}{【문제 1】}
A지방자치단체는 도로사업 부지를 취득하기 위하여 甲의 토지를 협의취득하여 공공용지의 협의취득을 원인으로 하는 소유권이전등기를 하였고, 乙의 토지에 대하여는 수용재결에 의하여 소유권을 취득한 후 소유권이전등기를 마쳤다. 그러나 甲과 乙의 토지(이하 ‘이 사건 토지’라 함)가 관내의 택지개발예정지구에 포함되자 A지방자치단체는 이 사건 토지가 도로사업에 더 이상 제공될 수 없는 상황에서 도로사업의 목적 달성이 불가능하다고 판단하여, 당초 협의취득 및 수용의 목적이 된 해당 도로사업을 폐지하였다. 이에 따라 甲과 乙에게 「공익사업을 위한 토지 등의 취득 및 보상에 관한 법률」에 의한 환매권이 발생하였다. 甲은 협의취득 당시에 수령한 보상금 상당 금액을 공탁한 후, A지방자치단체에게 환매의 의사표시를 하고 소유권이전등기청구소송을 제기하였다. 한편, 乙이 환매권을 행사할 무렵 환매금액에 관한 A지방자치단체와 乙의 협의가 성립되지 아니하여, A지방자치단체는 환매 대상 토지의 현재 가격이 취득일 당시에 비하여 현저히 상승하였음을 들어 환매대금의 증액을 구하는 소송을 제기하였다. 다음 물음에 답하시오. (40점) \par
\medskip\noindent 물음 1) 乙의 환매권 및 乙에 대한 환매대금증액청구소송의 법적 성질을 각각 설명하시오. (15점) \par
\medskip\noindent 물음 2) 甲의 소유권이전등기청구소송에서, A지방자치단체는 환매 대상 토지 가격의 상승에 따른 환매대금증액청구권을 내세워 증액된 환매대금과 보상금 상당액의 차액을 지급할 것을 선(先)이행 또는 동시이행의 항변으로 주장할 수 있는지에 관하여 설명하시오. (10점) \par
\medskip\noindent 물음 3) 만약 乙의 토지에 대한 수용재결에 취소사유에 해당하는 하자가 있어 乙이 환매권 행사 이전에 수용재결의 하자를 이유로 자신의 소유권 회복을 위한 소유권이전등기말소청구소송을 제기한 경우, 그 승소 여부를 검토하시오. (단, 수용재결에 불가쟁력이 발생하였음) (15점) \par
\medskip
\begin{framedtext}
\textbf{참조 조문} \par
\textbf{<공익사업을 위한 토지 등의 취득 및 보상에 관한 법률>} \par
\textbf{제91조(환매권)} ① 공익사업의 폐지·변경 또는 그 밖의 사유로 취득한 토지의 전부 또는 일부가 필요 없게 된 경우 토지의 협의취득일 또는 수용의 개시일(이하 이 조에서 “취득일”이라 한다) 당시의 토지소유자 또는 그 포괄승계인(이하 “환매권자”라 한다)은 다음 각 호의 구분에 따른 날부터 10년 이내에 그 토지에 대하여 받은 보상금에 상당하는 금액을 사업시행자에게 지급하고 그 토지를 환매할 수 있다. \par
\hspace*{1em}1. 사업의 폐지·변경으로 취득한 토지의 전부 또는 일부가 필요 없게 된 경우: 관계 법률에 따라 사업이 폐지·변경된 날 또는 제24조에 따른 사업의 폐지·변경 고시가 있는 날 \par
\hspace*{1em}2. 그 밖의 사유로 취득한 토지의 전부 또는 일부가 필요 없게 된 경우: 사업완료일 \par
②\textasciitilde{}③ <생략> \par
④ 토지의 가격이 취득일 당시에 비하여 현저히 변동된 경우 사업시행자와 환매권자는 환매금액에 대하여 서로 협의하되, 협의가 성립되지 아니하면 그 금액의 증감을 법원에 청구할 수 있다. \par
⑤ 제1항부터 제3항까지의 규정에 따른 환매권은 「부동산등기법」에서 정하는 바에 따라 공익사업에 필요한 토지의 협의취득 또는 수용의 등기가 되었을 때에는 제3자에게 대항할 수 있다. \par
\end{framedtext}
\end{longquestion}
